\section{Тензорная алгебра}

\subsection{Тензорные степени}

Пусть теперь~$V$ --- произвольный (возможно, свободный) $K$-модуль, а $V^*$ --- модуль, двойственный ему, то есть как множество состоящий из линейных функционалов:
\[
	V^* :=
	\Hom_{\Lin}(V, K) :=
	\set{\phi \colon V \to K}{\text{$\phi$ --- линейная}}
	.
\]

Будем писать
\[
	V^{\otimes m} :=
	\underbrace{V \oto V}_{m}
	,
\]
причём положим~$V^{\otimes 0} := K$, $V^{\otimes 1} := V$.

Тензорное произведение~$m$ экземпляров модуля~$V$ называется его \tdef{$m$-ой тензорной степенью}.

Рассмотрим прямую сумму
\[
	TV :=
	\bigoplus_{n \geq 0}
	V^{\otimes n}
	.
\]
В $n$-ом её слагаемом живут всевозможные тензоры <<длины~$n$>>
\[
	v_1 \oto v_n,
	\qquad
	v_i \in V,
\]
и их линейные комбинации. При этом элемент прямой суммы~$TV$ --- это \emph{конечная} сумма тензоров конечной длины, например
\[
	v_1 + w_1 \oto w_5 + u_1 \otimes u_2 \otimes u_3
	,
\]
где $v_1$, $w_i$, $u_j \in V$.

Тензорное умножение задаёт на~$TV$ структуру ассоциативной не обязательно коммутативной градуированной алгебры над~$K$.



Пусть $\ee = (e_1, \dots, e_n) \subset V$ --- базис в~$V$, а $\ee^* = (e^1, \dots, e^n) \subset V^*$ --- дуальный ему базис в~$V^*$. Напомню, что дуальность означает, что~$e^i(e_j) = \delta^i_j$.

Рассмотрим тензорное произведение вида
\[
	V^{\otp} \otimes (V^*)^{\otq}
	:=
	\underbrace{
		V \oto V
	}_{p}
	\otimes
	\underbrace{
		V^* \oto V^*
	}_{q}
	.
\]
Элементы
\[
	v_1 \oto v_p
	\otimes
	\xi^1 \oto \xi^q
	\in
	\Vtpq
\]
называются \tdef{разложимыми тензорами}. Если мы разложим каждый из сомножителей по базису, получим:
\begin{multline*}
	v_1 \oto v_p
	\otimes
	\xi^1 \oto \xi^q
	\nline{=}
	(v_1^{\attvec{i_1}}\,e_{\attvec{i_1}})
	\oto
	(v_p^{\attvec{i_p}}\,e_{\attvec{i_p}})
	\otimes
	(\xi^1_{\attcovec{j_1}}\,e^{\attcovec{j_1}})
	\oto
	(\xi^q_{\attcovec{j_q}}\,e^{\attcovec{j_q}})
	\nline{=}
	v_1^{\attvec{i_1}}
	\cdots
	v_p^{\attvec{i_p}}
	\cdot
	\xi^1_{\attcovec{j_1}}
	\cdots
	\xi^q_{\attcovec{j_q}}
	\cdot
	e_{\attvec{i_1}}
	\oto
	e_{\attvec{i_p}}
	\otimes
	e^{\attcovec{j_1}}
	\oto
	e^{\attcovec{j_q}}
	.
\end{multline*}
Множители 
$
	v_1^{\attvec{i_1}}
	\cdots
	v_p^{\attvec{i_p}}
	\cdot
	\xi^1_{\attcovec{j_1}}
	\cdots
	\xi^q_{\attcovec{j_q}}
$
при базисных тензорах называются \tdef{координатами} тензора. Так как произведение элементов кольца --- снова элемент кольца (например, произведение чисел --- это число), то можно составить произвольный тензор~$T \in \Vtpq$, задав его координатами в базисах~$(\ee, \ee^*)$:
\[
	T =
	T^{\attvec{i_1, \dots, i_p}}_{\attcovec{j_1, \dots, j_q}}
	\cdot
	e_{\attvec{i_1}} \oto e_{\attvec{i_p}}
	\otimes
	e^{\attcovec{j_1}} \oto e^{\attcovec{j_q}}
	.
\]
При этом 
$
	T^{\attvec{i_1, \dots, i_p}}_{\attcovec{j_1, \dots, j_q}}
	\in K
$ --- координаты тензора~$T$, а по наборам
\[
	(\attvec{i_1, \dots, i_p}) \in [n]^p
	\quad
	\text{и}
	\quad
	(\attcovec{j_1, \dots, j_q}) \in [n]^q
\]
предполагается суммирование.

\begin{remark}
	Обрати внимание, что координаты, относящиеся к \attvecem{векторам}, мы пишем \attvecem{сверху}, а координаты, относящиеся к \attcovecem{ковекторам}, пишем \attcovecem{снизу}.

	То, что относится к координатам \attvecem{векторов}, называется \attvecem{контравариантной} частью тензора, а то, что относится к координатам \attcovecem{ковекторов}, называется \attcovecem{ковариантной} его частью.

	Приставки \attcovecem{ко-} и \attvecem{контра-} описывают ситуацию с преобразованием координат: если $F \colon V \to V$ --- линейное отображение, которое переводит базис~$\ee$ в базис~$\hh$, то координаты \attcovecem{ковекторов} переводятся \attcovecem{в ту же сторону}, а координаты \attvecem{векторов} --- \attvecem{в обратную} (см.~\ref{eq:coords:diagram}).
\end{remark}

\begin{definition}
	Тензоры, живущие в~$\Vtpq$, называют \tdef{$p$ раз контравариантными и $q$ раз ковариантными} (по числу входящих координат векторов и ковекторов), а также \tdef{тензорами валентности~$\binom{p}{q}$} (по числу верхних и нижних индексов у координат).
\end{definition}

\subsubsection{Свёртка}

\begin{definition}
	Пусть $v_1 \oto v_p \otimes \xi^1 \oto \xi^q \in \Vtpq$. Его \tdef{свёрткой по паре индексов $(i, j) \in [p] \times [q]$} называется тензор
	\[
		\braket{\xi^j}{v_i}
		\cdot
		v_1 \oto \widehat{v_i} \oto v_p
		\otimes
		\xi^1 \oto \widehat{\xi^j} \oto \xi^q
		,
	\]
	в котором $j$-ый ковектор~$\xi^j$ подействовал на~$i$-ый вектор~$v_i$. Крышками обозначены сомножители, которых нет.
\end{definition}
Результирующий тензор будет иметь валентность~$\binom{p-1}{q-1}$.

Оператор свёртки по паре индексов~$(\attvec{i}, \attcovec{j}) \in \attvec{[p]} \times \attcovec{[q]}$ обозначим~$\conv{\attvec{i}}{\attcovec{j}}$, так как индекс~$\attvec{i}$ у~вектора пишется сверху, а индекс~$\attcovec{j}$ у~ковектора пишется сверху.

\begin{example}
	Найдём свёртку
	\[
		\conv{2}{1}
		\bigl(
			(e_1 + 3e_2 - 2e_3) \otimes
			(-e_1 + 2e_2) \otimes
			(e^1 + 3e^3) \otimes
			(e^4 - e^1)
		\bigr)
	\]
	по паре индексов~$(2, 1)$.

	Номер вектора равен~$2$, номер ковектора --- $1$. Значит, нужно свернуть второй вектор
	$
		(-e_1 + 2e_2)
	$
	и первый ковектор
	$
		(e^1 + 3e^3)
	$:
	\begin{align*}
		\braket{e^1 + 3e^3}{-e_1 + 2e_2} &=
		\braket{e^1}{-e_1} +
		\braket{e^1}{2e_2} +
		\braket{3e^3}{-e_1} +
		\braket{3e^3}{2e_2}
		\nlinea{=}
		-
		\braket{e^1}{e_1} +
		2\,
		\braket{e^1}{e_2}
		-3\,
		\braket{e^3}{e_1} +
		6\,
		\braket{e^3}{e_2}
		\nlinea{=}
		\delta_1^1 +
		2\,
		\delta^1_2
		-3\,
		\delta^3_1 +
		6\,
		\delta^3_2
		\nlinea{=}
		1
		.
	\end{align*}
	Тогда результат будет иметь вид
	\[
		1
		\cdot
		(e_1 + 3e_2 - 2e_3) \otimes
		(e^4 - e^1)
		.
	\]
\end{example}

